% Lab 2 --- compact version (full docs: docs/web/docs/labs/lab2/)
\section{Lab 2: Web Application Vulnerabilities --- SYNAPSE Intelligence Portal}

Lab~2 presents students with a custom two-tier web application,
\textit{SYNAPSE Intelligence Portal}, deployed on a segmented DMZ architecture.
The scenario is aligned with the OWASP Top~10~\cite{owasp2021} and chains
multiple vulnerabilities across two servers: Server-A hosts the portal front-end
with authentication and reporting features; Server-B hosts the DataVault internal
file management service accessible only after pivoting from Server-A. Students must
complete a progressive exploitation chain --- stored XSS, cookie theft, SQL injection,
and YAML deserialization --- to collect all three flags. A teacher reference --- credentials, startup procedure, and flag guide --- is compiled in Appendix~C;
complete source code, topology, and deployment instructions are at \path{docs/web/docs/labs/lab2/}.


\begin{figure}[H]
\centering
\includegraphics[width=0.95\textwidth]{lab2-topology.png}
\caption{Lab~2 network topology}
\label{fig:lab2-topology}
\end{figure}

\subsubsection*{Network Addressing}

\begin{table}[H]
\centering
\resizebox{\textwidth}{!}{\begin{tabular}{@{}lllll@{}}
\toprule
\textbf{Node} & \textbf{Role} & \textbf{Interface} & \textbf{IP} & \textbf{Network} \\
\midrule
Parrot-OS    & Attacker          & eth0  & 10.0.40.10       & 10.0.40.0/24 (Net-Attackers) \\
pfSense-LAB2 & Firewall/NAT      & WAN   & 192.168.0.x (DHCP) & 192.168.0.0/24 (Homelab)   \\
pfSense-LAB2 & Firewall/NAT      & LAN   & 172.16.1.1       & 172.16.1.0/30              \\
VyOS-LAB2    & DMZ Router        & eth0  & 172.16.1.2       & 172.16.1.0/30              \\
VyOS-LAB2    & DMZ Router        & eth1  & 192.168.30.1     & 192.168.30.0/24 (Net-DMZ)  \\
VyOS-LAB2    & DMZ Router        & eth2  & 10.0.40.1        & 10.0.40.0/24 (Net-Attackers)\\
Server-A     & SYNAPSE Portal    & eth0  & 192.168.30.10    & 192.168.30.0/24            \\
Server-B     & DataVault         & eth0  & 192.168.30.20    & 192.168.30.0/24            \\
\bottomrule
\end{tabular}}
\caption{Lab~2 IP addressing scheme}
\label{tab:lab2-ip}
\end{table}

\subsubsection*{OWASP Coverage and Vulnerability Chain}

\begin{table}[H]
\centering
\small
\begin{tabularx}{\textwidth}{llXl}
\toprule
\textbf{Step} & \textbf{OWASP} & \textbf{Vulnerability} & \textbf{Flag} \\
\midrule
1 & A03:2021 & Stored XSS in report field $\rightarrow$ session cookie theft & -- \\
2 & A07:2021 & Cookie injection $\rightarrow$ admin panel access & -- \\
3 & A03:2021 & SQL injection in search $\rightarrow$ credential dump & FLAG\{synapse\_sqli\_creds\_dumped\} \\
4 & A01:2021 & Lateral movement to Server-B via stolen creds & FLAG\{synapse\_pivot\_server\_b\} \\
5 & A08:2021 & YAML deserialization in DataVault upload $\rightarrow$ RCE & FLAG\{synapse\_rce\_achieved\} \\
\bottomrule
\end{tabularx}
\caption{Lab~2 OWASP Top~10 vulnerability chain}
\label{tab:lab2-vulns}
\end{table}

\subsubsection*{Results}

All three flags were successfully captured, confirming the complete exploitation
chain from stored XSS to remote code execution.

\begin{table}[H]
\centering
\resizebox{\textwidth}{!}{\begin{tabular}{@{}lll@{}}
\toprule
\textbf{Flag} & \textbf{Location} & \textbf{Technique} \\
\midrule
\texttt{FLAG\{synapse\_sqli\_creds\_dumped\}} & \texttt{/search?q=} UNION dump on Server-A & UNION SQL Injection \\
\texttt{FLAG\{synapse\_pivot\_server\_b\}}    & \texttt{/home/monitor/flag.txt} on Server-B & SSH pivot with SQLi-extracted credentials \\
\texttt{FLAG\{synapse\_rce\_achieved\}}       & \texttt{/preview} reverse shell on Server-B  & YAML deserialization RCE \\
\bottomrule
\end{tabular}}
\caption{Lab~2 flags and acquisition techniques}
\label{tab:lab2-flags}
\end{table}

\newpage

\subsection{Problems Encountered and Solutions}

\subsubsection{pfSense Block Private Networks and Reply-To}

Two pfSense default behaviours blocked traffic in the lab topology. First, the
\textit{Block private networks} option on the WAN interface was enabled by default,
causing pfSense to silently discard packets arriving from the simulated
attacker network (RFC~1918 space~\cite{rfc1918}). Second, the \textit{Reply-to}
feature added return routes in \texttt{pf} rules that conflicted with VyOS's
routing table, causing asymmetric routing. Both were disabled: the private network
block via the WAN interface settings, and reply-to via a per-rule option in the
firewall rule editor.

\subsubsection{VNC Keyboard Input on pfSense}

The same VNC key injection issue documented in Lab~1 reappeared on pfSense
console sessions, particularly when entering characters required by the pfSense
shell (\texttt{|}, \texttt{\&}, \texttt{\$}). The QEMU monitor raw socket approach
was applied identically.

\subsubsection{Server-A Volume Group Conflict}

When mounting the Server-A QCOW2 image from the EVE-NG host to pre-stage
application files, the host's own LVM volume group (\texttt{ubuntu-vg}) conflicted
with the guest disk's identically named group. Activating both simultaneously
caused the host's logical volumes to be shadowed. The solution was to rename the
guest volume group to \texttt{synapse-vg} using \texttt{vgrename} before activating
it on the host, then rename it back after dismounting.

\subsubsection{Server-B Installer Keyboard Layout}

During the Server-B installation, the Ubuntu installer defaulted to a Spanish
keyboard layout due to the locale of the EVE-NG host, causing mismatched characters
when entering the initial credentials. This was corrected by explicitly selecting the
\texttt{en\_US} keyboard layout at the installer prompt and verified with a
\texttt{/dev/urandom}-seeded password test before proceeding.

\subsubsection{Static Route Gateway Reference}

VyOS rejected a static route configuration that referenced a gateway IP not directly
reachable on any configured interface. The intended next-hop was on a subnet not yet
assigned to any VyOS interface at the time the route was applied. This was resolved
by defining the interface address before the static route, and by using
\texttt{commit} and \texttt{save} in the correct order within the VyOS configuration
session.

